\section*{الفصل \ref{ch:g9:metals-acids-corrosion} --- الفلزات: تفاعلاتها مع الأحماض وتآكلها}

\begin{solution}{exo:g9:metals-acids-corrosion:1}
الهيدروجين: عود مشتعل عند فوهة الأنبوب يُحدث فرقعة حادة.
\end{solution}

\begin{solution}{exo:g9:metals-acids-corrosion:2}
الحديد والزنك والألومنيوم تتفاعل؛ أما النحاس والذهب فلا.
\end{solution}

\begin{solution}{exo:g9:metals-acids-corrosion:3}
الأكسجين والماء معًا.
\end{solution}

\begin{solution}{exo:g9:metals-acids-corrosion:4}
\omterm{def:g9:ions:ion}{أيون} موجود في أثناء التفاعل لا يشارك فيه. ومع حمض الهيدروكلوريك: \omterm{def:g9:ions:ion}{أيون}
الكلوريد \ce{Cl-}.
\end{solution}

\begin{solution}{exo:g9:metals-acids-corrosion:5}
\ce{Mg(s) + 2H+(aq) -> Mg^{2+}(aq) + H2(g)}؛ والشحنة $+2$ في كل طرف.
\end{solution}

\begin{solution}{exo:g9:metals-acids-corrosion:6}
أضف هيدروكسيد الصوديوم: \omterm{def:g9:ions:precipitate}{فراسب} أخضر من \ce{Fe(OH)2} يبيّن \omterm{def:g9:ions:ion}{أيونات}
الحديد (II).
\end{solution}

\begin{solution}{exo:g9:metals-acids-corrosion:7}
الغليان يزيل \omterm{def:g4:air-a-mixture-of-gases:air}{الهواء} الذائب في الماء، والزيت يمنع \omterm{def:g4:air-a-mixture-of-gases:air}{الهواء} من أن \omterm{def:g2:mixing-and-dissolving:dissolve}{يذوب} فيه من
جديد: فيكون للمسمار ماء بلا أكسجين.
\end{solution}

\begin{solution}{exo:g9:metals-acids-corrosion:8}
الملح الذائب في الماء الذي يتناثر على السيارة يسرّع \omterm{def:g9:metals-acids-corrosion:corrosion}{الصدأ}.
\end{solution}

\begin{solution}{exo:g9:metals-acids-corrosion:9}
طبقة الأكسيد المتكوّنة تلتصق بالفلز وتعزله عن \omterm{def:g4:air-a-mixture-of-gases:air}{الهواء}: فيتوقف الهجوم عند
السطح.
\end{solution}

\begin{solution}{exo:g9:metals-acids-corrosion:10}
اطلِه، أو زيّته أو شحّمه، أو غلفنه (أو اصنعه من الفولاذ المقاوم \omterm{def:g9:metals-acids-corrosion:corrosion}{للصدأ}).
\end{solution}

\begin{solution}{exo:g9:metals-acids-corrosion:11}
\ce{2Al(s) + 6H+(aq) -> 2Al^{3+}(aq) + 3H2(g)}: 2 Al و6 H في كل طرف،
والشحنة $+6$ في كل طرف. ولكل 200 \omterm{def:g7:atoms-and-molecules:atom}{ذرة} Al: $200 \times \frac{3}{2} = 300$
\omterm{def:g7:atoms-and-molecules:molecule}{جزيء} هيدروجين.
\end{solution}

\begin{solution}{exo:g9:metals-acids-corrosion:12}
يتآكل الزنك المحيط بالخدش بدلًا من الفولاذ: فما دام في الجوار زنك يسلم
الفولاذ.
\end{solution}

\begin{solution}{pb:g9:metals-acids-corrosion:1}
\textbf{1.} \ce{Zn(s) + 2H+(aq) -> Zn^{2+}(aq) + H2(g)}.

\textbf{2.} الهيدروجين: فرقعة حادة مع عود مشتعل.

\textbf{3.} هيدروكسيد الصوديوم يعطي راسبًا أبيض من
\ce{Zn(OH)2}.

\textbf{4.} \omterm{def:g9:ions:ion}{أيونات} الكلوريد \ce{Cl-}.

\textbf{5.} $125.6 - 113.5 = \qty{12.1}{g}$.

\textbf{6.} ينطلق الهيدروجين ما دام الزنك يتفاعل مع الحمض؛ فإذا لم يبقَ
زنك توقف الفوران.

\textbf{7.} كان حديد الفولاذ سيتفاعل هو أيضًا، فيعطي الهيدروجين \omterm{def:g9:ions:ion}{وأيونات}
\ce{Fe^{2+}}؛ وكان هيدروكسيد الصوديوم سيعطي عندئذ راسبًا أخضر.

\textbf{8.} يُهاجَم الزنك بدلًا من الفولاذ، ويُبعد عنه الماء
\omterm{def:g4:air-a-mixture-of-gases:air}{والهواء}.

\textbf{9.} $2 \times 10.0 \times 10.0 = \qty{200}{cm^2}$.

\textbf{10.} $12.1 \div 7.13 \approx \qty{1.70}{cm^3}$.

\textbf{11.} $1.70 \div 200 = \qty{0.0085}{cm}$.

\textbf{12.} $0.0085 \times 10\,000 = \qty{85}{\micro m}$: فكان سُمك طبقة
الزنك نحو \qty{85}{\micro m}.
\end{solution}
